\scanpage{149}
\mspage{2}
% archive foliation: 148
II) Soit $a\in X$. ``Cône'' autour de $a$ :
\[
a\in C\subset X,\qquad
\underbrace{\underset{\substack{\Vert\\Y}}{\dot C}\unsure{\subset}C}_{\substack{\text{sous-espaces}\\\text{modérés}}}
\qquad\text{tels que}
\]
\[
(C,\unsure{a})\xleftarrow[\sim]{\varphi}
\deleted{\unsure{C_\varphi}}\ \operatorname{Cone}(\deleted{Y}\inserted{\dot C})
\qquad(\text{avec }\varphi(\unsure{a|o})=\unsure{o|a}).
\]
Cône strict : $\dot C$ est bordant dans $X\setminus(C\setminus\dot C)$.

Cône rigidifié.

\uline{Conj} a) \unsure{Existe} cône

b) L'espace des cônes \inserted{stricts} est contractile

c) L'espace des rigidifications des
cônes est contractile i.e. l'espace
des autom. de $\operatorname{Con}(Y)$, qui sont
l'identité sur $Y$ et fixent $a$,
est contractile (\unsure{Alexandroff} ?)

Donc le choix d'un cône strict autour
de $a$, ou d'un cône strict rigidifié,
ou d'une \deleted{cône} rigidification sur un
cône donné, est anodin.

$\overline{\underline{\mathrm{III}}}$

Catégorie \inserted{$PB$} des paires $(X,B)$\quad $B$ bordant dans $X$

Catégorie \inserted{$PBCR$} des \deleted{paires} triples
$(X,B,C,\deleted{\unsure{\dot C}}\varphi)$, des
paires bordantes avec collier \inserted{strict} rigidifié

Catégorie $PBC$\quad $(X,B,C,\dot C)$\quad (paires bord. avec coll.)

Catégorie \inserted{$EP$} des paires $(X,a)$\quad $X$ pointé
